\uuid{i9ZW}
\titre{Norme}
\auteur{bodin ; hanani ; mzari}
\organisation{exo7}
\licence{CC BY-NC-SA 4.0}
\datecreate{2020-01-20}
\niveau{L1}
\module{Géométrie}
\chapitre{Géométrie du plan}
\sousChapitre{Distance, norme, produit scalaire}
\difficulte{2}
\theme{norme}
\exo7qcmid{lille-1.geometrie-plane.014}

\contenu{
\begin{question}
\qclassification{140.01,140.02}
\qtags{motcle=norme}
Dans le plan muni d'une base orthonormée $(\vec{i},\vec{j})$, on considère deux vecteurs $\displaystyle \vec{u}$ et $\displaystyle \vec{v}$. On suppose que $\|\vec{u}\|=3$, $\|\vec{v}\|=3$ et que l'angle entre ces deux vecteurs est $\displaystyle \frac{\pi}{3}$. Quelle est la norme de $\vec{u}+\vec{v}$ ?
\begin{answers}
    \bad{$\displaystyle \|\vec{u}+\vec{v}\|=6$}
    \bad{$\displaystyle \|\vec{u}+\vec{v}\|=3$}
    \good{$\displaystyle \|\vec{u}+\vec{v}\|=3\sqrt{3}$}
    \bad{$\displaystyle \|\vec{u}+\vec{v}\|=9$}
\end{answers}
\begin{explanations}
La bilinéarité et la symétrie du produit scalaire donnent
$$\|\vec{u}+\vec{v}\|^2=(\vec{u}+\vec{v})\cdot(\vec{u}+\vec{v})=\|\vec{u}\|^2+\|\vec{v}\|^2+2\vec{u}\cdot\vec{v}.$$
Et puis $\displaystyle \vec{u}\cdot \vec{v}=\|\vec{u}\|\|\vec{v}\|\cos \left(\frac{\pi}{3}\right)=\frac{9}{2}$. Donc $\|\vec{u}+\vec{v}\|^2=9+9+9$.
\end{explanations}
\end{question}
}
